%% SCRAP: hardware/platform-integration/RASPBERRY_PI_BUILD
%% SOURCE: docs/working/hardware/platform-integration/RASPBERRY_PI_BUILD.adoc
%% STATUS: CURRENT
%% FITS: dev-guide/ch-raspi, user-guide/ch-raspi
%% EDITORIAL: lifted — prose rewritten to press voice

\section{Building StarForth on Raspberry Pi 4}

This guide covers building and tuning StarForth for the Raspberry Pi 4 with its
ARM64 assembly optimizations. The board pairs a Broadcom BCM2711 (quad-core
Cortex-A72 at 1.5\,GHz, ARMv8-A) with 32\,KB of L1 instruction and data cache
per core, 1\,MB of shared L2, and LPDDR4-3200 memory; NEON and CRC32 are both
present.

\subsection{Building}

The Pi 4 can be built natively or cross-compiled. A native build installs
\texttt{build-essential}, clones the repository, and compiles with the
Cortex-A72 tuning flags:

\begin{lstlisting}[language=bash]
make CFLAGS="-std=c99 -O3 -march=armv8-a+crc+simd -mtune=cortex-a72 \
  -DUSE_ASM_OPT=1 -DUSE_DIRECT_THREADING=1 -Iinclude \
  -Isrc/word_source -Isrc/test_runner/include"
\end{lstlisting}

Cross-compiling from x86\_64 Linux installs
\texttt{gcc-aarch64-linux-gnu}, builds with \texttt{CC=aarch64-linux-gnu-gcc}
and \texttt{LDFLAGS="-static"}, and copies the static binary to the board over
\texttt{scp}. The repository provides convenience targets: \texttt{rpi4} for a
native build, \texttt{rpi4-cross} for a static cross-build, and
\texttt{rpi4-perf} for the maximum-performance build with direct threading and
LTO.

\subsection{Architecture Detection and Code Integration}

A detection header (\texttt{include/arch\_detect.h}) selects the right
optimization headers and primitive macros at compile time --- x86\_64,
AArch64, or a portable fallback with a warning. Stack management, arithmetic
words, and dictionary lookup each keep an assembly path guarded by
\texttt{USE\_ASM\_OPT} alongside a C path, so the architecture-agnostic
\texttt{vm\_push\_opt}/\texttt{vm\_pop\_opt} macros resolve to the best
available implementation.

\subsection{Performance Tuning}

Three techniques dominate on the Cortex-A72. Aligning the \texttt{VM} struct to
the 64-byte cache line and grouping hot data (stacks and pointers) away from
cold data reduces false sharing. Prefetching the next dictionary entry while
comparing the current one hides traversal latency. And NEON block copy moves 64
bytes per iteration for bulk memory operations, falling back to scalar copy for
the remainder.

\subsection{Benchmarking}

Reliable measurement requires pinning the CPU to performance mode and watching
temperature. The governor is set with:

\begin{lstlisting}[language=bash]
echo performance | sudo tee /sys/devices/system/cpu/cpu*/cpufreq/scaling_governor
watch -n 1 vcgencmd measure_temp
vcgencmd get_throttled        # 0x0 means no throttling occurred
\end{lstlisting}

A benchmark script builds an \texttt{-O2} baseline and the optimized
\texttt{rpi4-perf} binary, runs a million iterations each of stack, arithmetic,
and logic loops, and reports the timing difference along with final temperature
and throttling status. The \texttt{perf} tool supplies cache-reference,
cache-miss, branch, and branch-miss counters for deeper analysis.

\subsection{Thermal and Power}

Sustained workloads need cooling. A passive aluminum or copper heatsink
covering the entire CPU with a 1\,mm thermal pad is the minimum; a 5\,V
30\,mm fan, optionally GPIO-controlled, handles continuous load. Power draw runs
2--3\,W idle, 4--5\,W under light load, and 7--8\,W with all cores busy, plus
2--3\,W for peripherals --- the official 5\,V/3\,A supply is recommended.
Overclocking through \texttt{/boot/config.txt} is possible but requires
adequate cooling and may void the warranty.

\subsection{Storage and Deployment}

For the SD card, choose UHS-I or better with an A2 application-performance
rating; brands with strong random I/O include Samsung EVO Plus, SanDisk
Extreme, and Kingston Canvas React. For production, a systemd service unit runs
StarForth on boot with automatic restart on failure; boot time can be trimmed
by disabling unused features and the splash screen in \texttt{/boot/config.txt};
and standard hardening (system updates, a UFW firewall, and key-only SSH)
rounds out a deployment.

\subsection{Troubleshooting}

\begin{itemize}
  \item \emph{Illegal instruction at runtime} --- the build exceeded the host's
    feature set; rebuild with the conservative \texttt{-march=armv8-a}.
  \item \emph{Cross-compiled binary will not run} --- confirm it is statically
    linked (\texttt{file build/starforth}); rebuild with \texttt{-static}
    otherwise.
  \item \emph{Slower than expected} --- verify the governor is
    \texttt{performance} rather than \texttt{ondemand} and check
    \texttt{vcgencmd get\_throttled} for thermal throttling.
  \item \emph{Out of memory} --- reduce \texttt{VM\_MEMORY\_SIZE} or add swap.
\end{itemize}
